% Contrastes pareados com a margem de equivalência (Seções 5.2.2 e 5.5.2).
% Dados: figuras/dados/contrastes_*.csv, gerados por
% `uv run python -m src.evaluation.export_contrastes`, com o mesmo cálculo das
% tabelas de contrastes. A faixa cinza é a margem; o veredito sai da posição do
% intervalo em relação a ela, e é isso que a figura deixa visível.
\begin{figure}[htb]
	\Caption{\label{fig:contrastes} Contrastes pareados e margens de equivalência}
	\centering
	\pgfplotsset{
		contraste/.style={
			width=0.62\textwidth, height=#1, clip=false,
			ytick=data, y dir=normal,
			axis lines*=left, tick align=outside,
			tick label style={font=\small}, label style={font=\small},
			xmajorgrids, grid style={draw=black!10},
			xticklabel style={/pgf/number format/use comma},
		},
		ponto/.style={only marks, mark=*, mark size=1.8pt, black,
			error bars/.cd, x dir=both, x explicit, error bar style={thick}},
		veredito/.style={only marks, mark=none, point meta=explicit symbolic,
			nodes near coords, nodes near coords align={right},
			every node near coord/.append style={font=\scriptsize\itshape, anchor=west}},
	}
	\pgfplotstableread[col sep=comma]{figuras/dados/contrastes_validade.csv}\tabvalidade
	\pgfplotstableread[col sep=comma]{figuras/dados/contrastes_fidelidade.csv}\tabfidelidade
	\begin{tikzpicture}
		\begin{axis}[contraste=3.4cm, name=validade,
				xmin=-25, xmax=25, ymin=0.4, ymax=2.6,
				xlabel={Diferença de validade (pontos percentuais)},
				yticklabels from table={\tabvalidade}{contraste},
				title={\small (a) Validade, margem de $\pm$10 pontos}]
			\fill[black!8] (axis cs:-10,0.4) rectangle (axis cs:10,2.6);
			\draw[dashed, black!50] (axis cs:0,0.4) -- (axis cs:0,2.6);
			\addplot[ponto] table[x=dif, y=y, x error minus=em, x error plus=ep, col sep=comma]
				{figuras/dados/contrastes_validade.csv};
			\addplot[veredito] table[x expr=26, y=y, meta=veredito, col sep=comma]
				{figuras/dados/contrastes_validade.csv};
		\end{axis}
		\begin{axis}[contraste=5cm, at={(validade.south west)}, anchor=north west, yshift=-2.5cm,
				xmin=-1.6, xmax=1.6, ymin=0.4, ymax=4.6,
				xlabel={Diferença de nota do juiz (escala de 0 a 5)},
				yticklabels from table={\tabfidelidade}{contraste},
				title={\small (b) Fidelidade (exploratória), margem de $\pm$0,5 ponto}]
			\fill[black!8] (axis cs:-0.5,0.4) rectangle (axis cs:0.5,4.6);
			\draw[dashed, black!50] (axis cs:0,0.4) -- (axis cs:0,4.6);
			\addplot[ponto] table[x=dif, y=y, x error minus=em, x error plus=ep, col sep=comma]
				{figuras/dados/contrastes_fidelidade.csv};
			\addplot[veredito] table[x expr=1.66, y=y, meta=veredito, col sep=comma]
				{figuras/dados/contrastes_fidelidade.csv};
		\end{axis}
	\end{tikzpicture}
	\Fonte{Elaborado pelo autor}
	\Nota{Ponto: diferença média pareada por item; barra: intervalo de confiança de 95\% por \textit{bootstrap}; faixa cinza: margem de equivalência. Intervalo inteiro dentro da faixa: equivalente; inteiro fora: diferente além da margem; nos demais casos, inconclusivo. A margem da validade foi fixada antes do cálculo dos intervalos; a da fidelidade, depois de vistas as notas.}
\end{figure}
